\begin{derivation}[从\eqnrefs{eq:cov-gauge-lagrangian-r2}到\eqnrefs{eq:covariant-photon-mode-comm-r7}]
在 Feynman 规范 $\xi_{\rm gf}=1$ 中，先把
\[
\mathcal L_{\rm EM,1}
=-\frac1{4\mu_0}F_{\mu\nu}F^{\mu\nu}
-\frac1{2\mu_0}(\partial_\mu A^\mu)^2
\]
展开。利用
\[
F_{\mu\nu}F^{\mu\nu}
=2\bigl[(\partial_\mu A_\nu)(\partial^\mu A^\nu)
-(\partial_\mu A_\nu)(\partial^\nu A^\mu)\bigr]
\]
并对第二项分部积分，交叉项与规范固定项相消；忽略不影响 Euler--Lagrange 方程的全微分后，二次拉氏量变成
\[
\mathcal L_{\rm EM,1}
=-\frac1{2\mu_0}(\partial_\mu A_\nu)(\partial^\mu A^\nu).
\]
因此以 $A^\nu$ 为正则坐标时
\[
\Pi_\nu
\equiv\frac{\partial\mathcal L_{\rm EM,1}}
{\partial(\partial_t A^\nu)}
=-\frac1{\mu_0c^2}\partial_tA_\nu
=-\varepsilon_0\partial_tA_\nu.
\]
这里的下指标十分重要：Minkowski 度规已经进入 $A_\nu=\eta_{\nu\rho}A^\rho$，所以四个分量不会得到同号的正定振子度量。

令
\[
N_{\boldsymbol\kappa}
\equiv\sqrt{\frac{\hbar}{2\varepsilon_0\omega_{\boldsymbol\kappa}}}
\]
并把正文\eqnrefs{eq:covariant-photon-fourpol-r7} 写成
\[
A^\mu(x)=\sum_\lambda\int\frac{\dd^3\kappa}{(2\pi)^3}N_{\boldsymbol\kappa}
\left(\epsilon_\lambda^\mu a_\lambda \ee^{-\ii\kappa\cdot x}
+\epsilon_\lambda^{\mu *}a_\lambda^\dagger \ee^{+\ii\kappa\cdot x}\right).
\]
时间微分给出 $\mp\ii\omega_{\boldsymbol\kappa}$，所以
\[
\Pi_\nu(x)=\ii\varepsilon_0\sum_\lambda\int\frac{\dd^3\kappa}{(2\pi)^3}
\omega_{\boldsymbol\kappa}N_{\boldsymbol\kappa}
\left(\epsilon_{\lambda\nu}a_\lambda \ee^{-\ii\kappa\cdot x}
-\epsilon_{\lambda\nu}^*a_\lambda^\dagger \ee^{+\ii\kappa\cdot x}\right).
\]
暂不预设模代数，而写
\[
[a_\lambda(\boldsymbol\kappa),a_{\lambda'}^\dagger(\boldsymbol\kappa')]
=C_{\lambda\lambda'}(\boldsymbol\kappa)
(2\pi)^3\delta^{(3)}(\boldsymbol\kappa-\boldsymbol\kappa').
\]
在等时 $x^0=y^0$ 计算 $[A^\mu(\mathbf x),\Pi_\nu(\mathbf y)]$。只有 $a$--$a^\dagger$ 与 $a^\dagger$--$a$ 两类收缩；用 $N_{\boldsymbol\kappa}^2\varepsilon_0\omega_{\boldsymbol\kappa}=\hbar/2$ 后，两类正负频贡献恰好合成
\begin{align*}
[A^\mu(\mathbf x),\Pi_\nu(\mathbf y)]
&=-\ii\hbar\int\frac{\dd^3\kappa}{(2\pi)^3}
 \ee^{\ii\boldsymbol\kappa\cdot(\mathbf x-\mathbf y)}
 \sum_{\lambda\lambda'}C_{\lambda\lambda'}
 \epsilon_\lambda^\mu\epsilon_{\lambda'\nu}^*.
\end{align*}
正则量子化要求这等于
\[
\ii\hbar\delta^\mu{}_{\nu}\delta^{(3)}(\mathbf x-\mathbf y)
=\ii\hbar\int\frac{\dd^3\kappa}{(2\pi)^3}
 \ee^{\ii\boldsymbol\kappa\cdot(\mathbf x-\mathbf y)}\delta^\mu{}_{\nu}.
\]
由于四偏振基满足正文\eqnrefs{eq:fourpol-completeness-r7}，即
\[
\sum_{\lambda\lambda'}\eta^{\lambda\lambda'}
\epsilon_\lambda^\mu\epsilon_{\lambda'\nu}^*=\delta^\mu{}_{\nu},
\]
逐个 Fourier 模比较得到唯一解
\[
\boxed{C_{\lambda\lambda'}=-\eta_{\lambda\lambda'}}.
\]
于是
\[
[a_\lambda(\boldsymbol\kappa),a_{\lambda'}^\dagger(\boldsymbol\kappa')]
=-\eta_{\lambda\lambda'}(2\pi)^3\delta^{(3)}(\boldsymbol\kappa-\boldsymbol\kappa'),
\]
即正文\eqnrefs{eq:covariant-photon-mode-comm-r7}。特别地，在这个不定内积空间中，单个时间样激发的自内积为
\[
\langle0|a_0a_0^\dagger|0\rangle<0,
\]
所以负范数并非另加的假设，而是“显式 Lorentz 协变 + 四分量正则量子化 + Minkowski 度规”三者共同固定的代数结果。Gupta--Bleuler 条件随后才负责从这个不定内积空间中选出物理商空间。
\end{derivation}
\begin{derivation}[从\eqnrefs{eq:physical-photon-polarization-constraints-r68}到\eqnrefs{eq:physical-photon-polarization-projector-dp42}]
取实光子的物理四动量 $k^\mu$，满足 $k^2=0$，并选一个不与 $k$ 正交的辅助四矢量 $n^\mu$。物理偏振 $\epsilon_\lambda^\mu$ 满足

\begin{equation*}
 k\!\cdot\epsilon_\lambda=0,
 \qquad
 n\!\cdot\epsilon_\lambda=0,
 \qquad
 \epsilon_\lambda\!\cdot\epsilon_{\lambda'}^*=-\delta_{\lambda\lambda'},
 \qquad \lambda=1,2.
\end{equation*}
偏振和
\begin{equation*}
 P^{\mu\nu}(k,n)
 \equiv\sum_{\lambda=1,2}
 \epsilon_\lambda^\mu\epsilon_\lambda^{\nu*}
\end{equation*}
必须由 $\eta^{\mu\nu},k^\mu,n^\mu$ 构成，并满足
\begin{equation*}
 P^{\mu\nu}k_\nu=0,
 \qquad
 P^{\mu\nu}n_\nu=0.
\end{equation*}
从一般张量设式
\begin{equation*}
 P^{\mu\nu}
 =-\eta^{\mu\nu}
 +A(k^\mu n^\nu+n^\mu k^\nu)
 +B k^\mu k^\nu
 +C n^\mu n^\nu
\end{equation*}
出发。与 $k_\nu$ 收缩并用 $k^2=0$：
\begin{equation*}
 0=-k^\mu+A(k\!\cdot n)k^\mu+C(k\!\cdot n)n^\mu.
\end{equation*}
由于 $k^\mu,n^\mu$ 独立，
\begin{equation*}
 A=\frac1{k\!\cdot n},
 \qquad C=0.
\end{equation*}
再与 $n_\nu$ 收缩：
\begin{equation*}
 0=-n^\mu+\frac{k^\mu n^2+n^\mu(k\!\cdot n)}{k\!\cdot n}
 +B(k\!\cdot n)k^\mu,
\end{equation*}
从 $k^\mu$ 系数得到
\begin{equation*}
 B=-\frac{n^2}{(k\!\cdot n)^2}.
\end{equation*}
因此
\begin{equation*}
 {
 P^{\mu\nu}(k,n)
 =-\eta^{\mu\nu}
 +\frac{k^\mu n^\nu+n^\mu k^\nu}{k\!\cdot n}
 -\frac{n^2k^\mu k^\nu}{(k\!\cdot n)^2}.}
 \end{equation*}
若振幅的光子指标与守恒流 $J_\mu$ 收缩，并满足 $k^\mu J_\mu=0$，则
\begin{equation*}
 {
 J_\mu P^{\mu\nu}(k,n)J_\nu^*
 =-J_\mu J^{\mu*}.}
 \end{equation*}
所有依赖辅助矢量 $n^\mu$ 的项都消失。因此在规范不变可观测量中，物理偏振求和常可在与守恒振幅收缩之后简写为 $-\eta^{\mu\nu}$；这不是偏振空间本身的恒等式。
\end{derivation}

\begin{derivation}[从\eqnrefs{eq:dirac-field-pole-overlap-dp8}到\eqnrefs{eq:dirac-lsz-amputation-r68}]

设任意连通 Green 函数含一条出射 Dirac 外腿，写成
\begin{equation*}
 \widetilde G_\alpha(p;\{q\})
 =\int\dd^4 x\,\ee^{+\ii p\cdot x/\hbar}
 \langle\Omega|T\{\psi_\alpha(x)\mathcal O(\{q\})\}|\Omega\rangle_c.
\end{equation*}
若完整谱在 $p^2=m_{\rm phys}^2c^2$ 处含稳定单粒子极点，则极点邻域必可因子化为
\begin{equation*}
 \widetilde G(p;\{q\})
 =\widetilde S(p)\,\widetilde\Gamma(p;\{q\})
 +\widetilde G_{\rm reg}(p;\{q\}),
\end{equation*}
其中 $\widetilde\Gamma$ 与 $\widetilde G_{\rm reg}$ 在该极点解析。正文\eqnrefs{eq:dirac-field-pole-overlap-dp8} 给出场--粒子重叠
\begin{equation*}
 \langle\Omega|\psi(0)|p,s\rangle=\sqrt{Z_{\rm pole}^{(e)}}\,u_s(p),
\end{equation*}
所以精确二点函数的单粒子极点部分为
\begin{equation*}
 \widetilde S(p)
 \xrightarrow[p^2\to m_{\rm phys}^2c^2]{}
 \frac{\ii\hbar Z_{\rm pole}^{(e)}(\slashed p+m_{\rm phys}c)}
 {p^2-m_{\rm phys}^2c^2+\ii0^+}
 +\widetilde S_{\rm reg}(p).
\end{equation*}
左乘自由逆 Dirac 核并使用
\begin{equation*}
 (\slashed p-mc)(\slashed p+mc)=p^2-m^2c^2
\end{equation*}
可知
\begin{equation*}
 [Z_{\rm pole}^{(e)}]^{-1/2}\frac{\slashed p-mc}{\ii\hbar}\,\widetilde S(p)
 \xrightarrow[p^2\to m^2c^2]{}\sqrt{Z_{\rm pole}^{(e)}}
\end{equation*}
在取单粒子极点留数的意义下成立。于是
\begin{equation*}
 \lim_{p^2\to m^2c^2}
 \bar u_s(p)\,
 [Z_{\rm pole}^{(e)}]^{-1/2}\frac{\slashed p-mc}{\ii\hbar}
 \widetilde G(p;\{q\})
 =\bar u_s(p)\widetilde\Gamma(p;\{q\}),
\end{equation*}
即正文\eqnrefs{eq:dirac-lsz-amputation-r68}。连续谱和其它在该点解析的部分没有相同极点，因此在取留数时消失；严格无质量 QED 的红外粒子结构需按正文给出的调节/包容方案另行处理。
\end{derivation}

\begin{derivation}[从\eqnrefs{eq:photon-physical-momentum-inverse-kernel-dp42}到\eqnrefs{eq:photon-lsz-amputation-r68}]
物理单光子态 $|k,\lambda\rangle$ 只含两个横向偏振，$\lambda=1,2$。其场--粒子重叠可写成

\begin{equation*}
 \langle\Omega|A_\mu(0)|k,\lambda\rangle
 =\sqrt{Z_{\rm pole}^{(\gamma)}}\,\varepsilon_\mu^{(\lambda)}(k),
 \qquad
 k^2=0,
 \qquad
 k\cdot\varepsilon^{(\lambda)}=0.
\end{equation*}
因此精确光子二点函数在横向单光子极点附近具有
\begin{equation*}
 \widetilde D_{\mu\nu}(k)
 =Z_{\rm pole}^{(\gamma)}\,\widetilde D^{(0)}_{\mu\nu}(k)
 +\widetilde D^{\rm reg}_{\mu\nu}(k),
\end{equation*}
其中 $\widetilde D^{\rm reg}_{\mu\nu}$ 在该极点邻域解析或属于连续谱部分。物理四动量表示的自由逆 Maxwell 核由\eqnrefs{eq:photon-physical-momentum-inverse-kernel-dp42} 定义。结合\eqnrefs{eq:photon-prop-kappa-to-q-r13}，可直接相乘得到
\begin{align*}
 \mathscr K_{\xi_{\rm gf}}^{\mu\rho}(k)
 \widetilde D^{(0)}_{\rho\nu}(k)
 &=\left[-\frac{k^2}{\mu_0\hbar^2}
 \left(P_T^{\mu\rho}+\frac1{\xi_{\rm gf}}P_L^{\mu\rho}\right)\right]\\
 &\quad\times
 \left[-\frac{\ii\mu_0\hbar^3}{k^2+\ii0^+}
 \left(P_{T,\rho\nu}+\xi_{\rm gf}P_{L,\rho\nu}\right)\right],\\
 &\xrightarrow[k^2\to0]{}\ii\hbar\,
 (P_T+P_L)^\mu{}_{\nu},\\
 &=\ii\hbar\,\delta^\mu{}_{\nu}.
\end{align*}
第二行使用 $P_T^2=P_T$、$P_L^2=P_L$ 与 $P_TP_L=0$。于是精确传播子的极点部分满足
\begin{equation*}
 [Z_{\rm pole}^{(\gamma)}]^{-1/2}\frac{\mathscr K_{\xi_{\rm gf}}^{\mu\rho}(k)}{\ii\hbar}
 \widetilde D_{\rho\nu}(k)
 \xrightarrow[k^2\to0]{}
 [Z_{\rm pole}^{(\gamma)}]^{1/2}\delta^\mu{}_{\nu}.
\end{equation*}
再与 $\varepsilon_\mu^{(\lambda)*}$ 或 $\varepsilon_\mu^{(\lambda)}$ 收缩，便得到物理光子外态。因而 $(\ii\hbar)^{-1}$ 不是可省略的约定因子，而是由二点函数的归一化由前述约定固定。

不同协变规范只改变含 $k_\mu$ 或 $k_\nu$ 的纵向部分。物理偏振满足
\begin{equation*}
 k_\mu\varepsilon^{\mu(\lambda)}=0,
\end{equation*}
而与守恒物质流连接时 Ward 恒等式给出
\begin{equation*}
 k_\mu\mathcal M^\mu=0.
\end{equation*}
因此 $\varepsilon^\mu\mapsto\varepsilon^\mu+\alpha k^\mu$ 不改变散射矩阵元；光子外腿表示规范等价类中的横向物理态，而不是四势的四个分量各自对应一个粒子。
\end{derivation}

\begin{derivation}[从\eqnrefs{eq:lsz-dirac-photon-mainline-dp8}到\eqnrefs{eq:multi-leg-lsz-r68}]
考虑含 $N_f$ 条费米子外腿和 $N_\gamma$ 条光子外腿的连通精确 Green 函数。若每一个外部不变量都分别接近一个孤立单粒子极点，多变量 Laurent 展开的最高阶极点部分必可分解为

\begin{equation*}
 \widetilde G_{\rm conn}
 =\left(\prod_{a=1}^{N_f}\widetilde S_a\right)
 \left(\prod_{b=1}^{N_\gamma}\widetilde D_b\right)
 \widetilde\Gamma_{\rm amp}
 +\text{至少缺少一个外腿极点的项}.
\end{equation*}
不同外部坐标具有独立 Fourier 共轭动量，因此逐个外部不变量取留数等价于逐条外腿提取局部 Laurent 系数。对每条费米子外腿乘
\begin{equation*}
 [Z_{\rm pole}^{(e)}]^{-1/2}\frac{\slashed p-mc}{\ii\hbar},
\end{equation*}
对每条光子外腿乘
\begin{equation*}
 [Z_{\rm pole}^{(\gamma)}]^{-1/2}\frac{\mathscr K(k)}{\ii\hbar},
\end{equation*}
再投影到相应壳上自旋子和横向偏振后，所有外腿传播极点及其留数均被除去，得到
\begin{equation*}
 \mathcal M_{fi}
 =\left[\prod_{\text{外部费米子}}\text{壳上 }u,\bar u,v,\bar v\right]
 \left[\prod_{\text{外部光子}}\varepsilon^{(\lambda)}\right]
 \widetilde\Gamma_{\rm amp}.
\end{equation*}
所有相互作用顶点的平移不变性共同产生整体
$(2\pi\hbar)^4\delta^{(4)}(P_f-P_i)$；抽去这一总四动量守恒因子后，$\mathcal M_{fi}$ 才进入截面和衰变率公式。
\end{derivation}
